% ECONOMICS_PROBLEM: {"id": "C73-30", "title": "Weak master-equation selection beyond finite-state potentials", "jel_code": "C73", "source_id": "OP-0216"}
% !TeX program = xelatex
\documentclass[11pt,a4paper]{article}
\usepackage[margin=23mm]{geometry}
\usepackage{fontspec}
\setmainfont{Latin Modern Roman}
\usepackage{amsmath,amssymb,mathtools}
\usepackage{xcolor,enumitem,booktabs,longtable,array,xurl,hyperref}
\definecolor{muted}{HTML}{53636C}
\hypersetup{colorlinks=true,urlcolor=blue,linkcolor=blue}
\setlength{\parindent}{0pt}
\setlength{\parskip}{5pt}
\newcommand{\R}{\mathbb R}
\newcommand{\E}{\mathbb E}
\newcommand{\N}{\mathbb N}
\newcommand{\ind}{\mathbf 1}
\newcommand{\Var}{\operatorname{Var}}
\newcommand{\Cov}{\operatorname{Cov}}
\newcommand{\supp}{\operatorname{supp}}
\DeclareMathOperator*{\argmin}{arg\,min}
\DeclareMathOperator*{\argmax}{arg\,max}
\newcommand{\entrymeta}[3]{{\sffamily\small\color{muted}\textbf{\texttt{#1}}\quad|\quad #2\par #3\par}}
\newcommand{\Problemref}[1]{\texttt{#1}}
\newcommand{\JELref}[2]{\texttt{#2} (JEL \texttt{#1})}
\begin{document}
\section*{Weak master-equation selection beyond finite-state potentials}
\textbf{Problem ID:} \texttt{C73-30}\quad \textbf{JEL:} \texttt{C73}\par
% BEGIN ECONOMICS STATEMENT
\label{problem:OP-0216}
% GAME_THEORY_PROBLEM: {"local_id": "GT-086", "original_number": 86, "kind": "R", "entry_kind": "statement_only", "published": false, "review_required": true, "category": "game_theory"}
\label{game:GT-086}
{\sffamily\small\color{muted}
{\small\textbf{GT-086} | Mean-field games | R: corrected nonpotential extension; the potential uniqueness subcase is known\par}
\par}
\subsection*{Definitions and mathematical statement}
Let $E=\{1,\ldots,d\}$, $d\geq3$, $\Delta_d=\{m\geq0:\sum_i m_i=1\}$, $T>0$, and smooth bounded $f,g:\Delta_d\to\mathbb R^d$, with no potential or monotonicity assumption. Set
\[
 H_i(U)=\tfrac12\sum_{j\ne i}(U_i-U_j)_+^2,\qquad
 b_i(m,U)=\sum_{j\ne i}[m_j(U_j-U_i)_+-m_i(U_i-U_j)_+].
\]
The deterministic classical master equation is
\[
 -\partial_tU_i+H_i(U)-\sum_j b_j(m,U)\partial_{m_j}U_i=f_i(m),
 \qquad U_i(T,m)=g_i(m).
\]
For nonsmooth $U$, products involving $b(m,U)DU$ do not acquire a distributional meaning merely by writing this display.

Here is an explicit selection formulation of the agenda. Let $\mathcal A(f,g,T)$ be the class of smooth nonnegative boundary-repulsion families $\phi_\varepsilon$ tending to zero on each compact subset of $(0,1]$ for which the following parabolic problem has a unique classical solution. Require that solution to be continuous on $[0,T]\times\Delta_d$, $C^1$ in time and $C^2$ in the interior, with bounded value and first derivatives and second derivatives bounded on interior compact sets. Put
\[
 b_i^\phi=b_i+\sum_{j\ne i}[m_j\phi_\varepsilon(m_i)-m_i\phi_\varepsilon(m_j)],\quad
 L_\varepsilon=\varepsilon\sum_{j,k}(m_j\delta_{jk}-m_jm_k)\partial_{m_jm_k}^2.
\]
The regularized field $U^{\varepsilon,\phi}$ satisfies
\[
 \begin{split}
 0={}&\partial_tU_i+L_\varepsilon U_i+\sum_j b_j^\phi(m,U)\partial_{m_j}U_i
 +2\varepsilon\sum_jm_j(\partial_{m_i}U_i-\partial_{m_j}U_i)\\
 &-H_i(U)+f_i(m)+\sum_{j\ne i}\phi_\varepsilon(m_j)(U_j-U_i),
 \qquad U_i(T,m)=g_i(m).
 \end{split}
\]
Tangent derivatives make the equation independent of extensions off the simplex. Let $Z^{\varepsilon,\phi}=(U_i^{\varepsilon,\phi}-U_d^{\varepsilon,\phi})_{i<d}$. These are the differences determining equilibrium rates. Define a \emph{selected weak feedback field} as an $L^1_{\mathrm{loc}}((0,T)\times\operatorname{int}\Delta_d)$ subsequential limit of these $Z$ fields, where Lebesgue measure is the $(d-1)$-dimensional simplex measure.

Find assumptions, genuinely beyond the potential class, and an intrinsic admissibility condition under which: (i) such limits exist for admissible regularizations; (ii) every admissible selected field is the same almost everywhere; and (iii) it agrees with the component differences of the classical master solution whenever that solution exists. One precise independence target is
\[
 \lim_{\varepsilon\downarrow0}
 \|Z^{\varepsilon,\phi}-Z^{\varepsilon,\psi}\|_{L^1(K)}=0
 \quad\text{for all }\phi,\psi\in\mathcal A(f,g,T)
 \text{ and compact }K\subset(0,T)\times\operatorname{int}\Delta_d,
\]
together with compactness. Different cost corrections or boundary conventions define different admissible classes and must be named.
\subsection*{Short English statement}
Can a well-defined weak feedback solution and uniqueness criterion be obtained for nonpotential finite-state master equations after common noise disappears?
\subsection*{Variants and scope boundaries}
Feedback differences have dimension $d-1$; the full value vector also contains a common component. Uniqueness among arbitrary distributions is a different, unjustified claim. Potentiality provides a conservative formulation and a weak one-sided-Lipschitz criterion in the cited work; that subcase is already known.
\subsection*{Source relationship and status}
The input's blanket finite-state uniqueness description needs correction: the chapter supplies a potential weak uniqueness result. Its nonpotential extension motivates the editorial agenda above.
\subsection*{References}
\begingroup\footnotesize\raggedright
{\footnotesize\raggedright
\textbf{[S1]} Cardaliaguet and Delarue. \emph{Selected topics in mean field games}, ICM2022, vol.~5; DOI: 10.4171/ICM2022/17.\par
\textit{Locators:} Section~4.3, p.~3689, loss of classical solutions and weak uniqueness issue; Theorem~4.6 and following conservative-form paragraph, p.~3692, known potential solution; Section~4.4, p.~3692, nonpotential extension.\par
\url{https://ems.press/content/book-chapter-files/33265}\par\smallskip
\textbf{[S2]} Alekos Cecchin and Fran\c{c}ois Delarue. \emph{Selection by vanishing common noise for potential finite state mean field games}. Communications in Partial Differential Equations 47(1) (2022), 89--168; DOI: 10.1080/03605302.2021.1955256.\par
\textit{Locator/source role:} abstract's final intrinsic uniqueness criterion, and the chapter's Theorem~4.6 reference [43]; this is a known boundary, not the source of a universal nonpotential theorem.\par
\url{https://arxiv.org/abs/2005.12153}\par}
\par\endgroup

\subsection*{Common conventions: Source mathematical conventions}

\textbf{Finite games.} For a finite set $A$, $\Delta(A)$ is its probability
simplex. Players choose independent mixed strategies in Nash equilibrium;
correlation is allowed only when explicitly part of the solution concept.
In a game with payoff functions $u_i$ and strategy profile $x$, an additive
$\varepsilon$ Nash equilibrium satisfies
\[
 u_i(x)\geq u_i(x_i',x_{-i})-\varepsilon
 \quad\text{for every player $i$ and every admissible deviation $x_i'$}.
\]
For numerical approximation, payoffs are normalized as locally specified.
An $\varepsilon$ payoff residual is distinct from distance at most
$\varepsilon$ to the set of exact equilibria. Point convergence, convergence
of empirical play, and convergence in distance to an equilibrium set are
also distinct.

\textbf{Computational representations.} Unless stated otherwise, a complexity
question takes rational data with binary encoding; polynomial time is measured
in total bit length. A fixed-accuracy polynomial algorithm is not necessarily
polynomial in $1/\varepsilon$, and neither is necessarily polynomial in
$\log(1/\varepsilon)$. Succinct games and value, demand, or sampling oracles
must declare their interfaces. Exact real solutions require an explicit
representation; an unspecified list of arbitrary real numbers is not a
finite computational output. A claimed algorithmic barrier is meaningful
only for its specified model and reductions.

\textbf{Dynamic games.} Histories, information, observation, and allowable
randomization are part of the model. A stationary strategy depends only on
the current state; a positional strategy in a deterministic graph game
chooses one action at each controlled vertex. Nash equilibrium at an initial
state and equilibrium after every history are different requirements.
An expectation of a pathwise $\liminf$ payoff, a $\liminf$ of expected averages,
a limiting expected average, and a uniform finite-horizon equilibrium are
different criteria. None is silently substituted for another.

\textbf{Allocation and mechanisms.} A complete allocation assigns every item
exactly once unless otherwise stated. Zero-valued goods, negative values,
capacity constraints, feasibility, lotteries, and copies of goods affect
fairness definitions and are specified locally. Dominant-strategy incentive
compatibility, Bayesian incentive compatibility, universal truthfulness,
and truthfulness in expectation impose different inequalities. Whenever
an objective ratio has zero denominator, its convention is stated or those
instances are excluded.

\textbf{Infinite-dimensional models.} Expectations, integrals, stochastic
equations, and optimization objectives require their local regularity,
integrability, filtrations, and admissible domains. A backward--forward
mean-field system is a boundary-value problem; it is not an initial-value
semiflow without an additional construction. Long-time, large-population,
and vanishing-noise limits require an order of limits and a convergence
topology. If a minimum need not be attained, the objective is its infimum
or supremum and the approximation requirement is stated separately.
% END ECONOMICS STATEMENT
\end{document}
